% !TeX program = tectonic
\documentclass[11pt,a4paper]{article}
\usepackage{fontspec}
\usepackage[russian]{babel}
\babelfont{rm}[Language=Default]{Liberation Serif}
\babelfont[Language=Default]{sf}{Carlito}
\babelfont[Language=Default]{tt}{DejaVu Sans Mono}
\usepackage{amsmath,amssymb,amsthm}
\usepackage{geometry}
\geometry{a4paper, top=2.4cm, bottom=2.4cm, left=2.4cm, right=2.4cm}
\usepackage{booktabs,tabularx,enumitem,xcolor,tcolorbox}
\tcbuselibrary{breakable,skins}
\definecolor{linkblue}{HTML}{1F4E79}
\definecolor{statgreen}{HTML}{2F6B3A}
\definecolor{goldacc}{HTML}{8B7E5A}
\usepackage[colorlinks=true,linkcolor=linkblue,urlcolor=linkblue,unicode=true]{hyperref}
\theoremstyle{plain}
\newtheorem{theorem}{Теорема}
\newtheorem{lemma}[theorem]{Лемма}
\newtheorem{proposition}[theorem]{Предложение}
\newtheorem{corollary}[theorem]{Следствие}
\theoremstyle{definition}
\newtheorem{definition}[theorem]{Определение}
\theoremstyle{remark}
\newtheorem{remark}[theorem]{Замечание}
\newtcolorbox{statusbox}[1][]{colback=statgreen!5,colframe=statgreen!75,
  title={\textbf{Сводка доказанного:} #1},fonttitle=\small,boxrule=0.7pt,arc=2pt,breakable}
\newtcolorbox{databox}[1][]{colback=goldacc!6,colframe=goldacc!80,
  title={\textbf{Данные протокола:} #1},fonttitle=\small,boxrule=0.7pt,arc=2pt,breakable}
\newtcolorbox{verifybox}[1][]{colback=linkblue!5,colframe=linkblue!80,
  title={\textbf{Верификация:} #1},fonttitle=\small,boxrule=0.7pt,arc=2pt,breakable}
\widowpenalty=10000
\clubpenalty=10000
\setlength{\parskip}{0.35em}
\newcommand{\CC}{\mathbb{C}}
\newcommand{\RR}{\mathbb{R}}
\newcommand{\ZZ}{\mathbb{Z}}
\newcommand{\QQ}{\mathbb{Q}}
\newcommand{\Ga}{\Gamma}
\newcommand{\Bfun}{\mathrm{B}}
\newcommand{\im}{\operatorname{Im}}
\newcommand{\re}{\operatorname{Re}}
\newcommand{\lcm}{\operatorname{lcm}}
\newcommand{\SNF}{\operatorname{SNF}}
\newcommand{\Jac}{\operatorname{Jac}}
\newcommand{\rank}{\operatorname{rank}}
\newcommand{\Ind}{\operatorname{Index}}
\newcommand{\disc}{\operatorname{disc}}
\begin{document}
\begin{center}
{\LARGE\bfseries Γ-тождества}\\[0.4em]
{\large формула отражения и формула умножения Гаусса}\\[0.6em]
{\normalsize Исаев Исхак Хамзатович}\\[0.2em]
{\small Программа <<Динамический принцип>> --- лаборатория \texttt{hodge-laboratory}}\\[0.15em]
{\small Отдельная монография теоремы · Монография 3/16}
\end{center}
\vspace{0.4em}
{\color{linkblue}\hrule height 0.8pt}
\vspace{0.8em}

\section{Постановка}

Весь аналитический инструментарий стенда $N=15/30$ --- два
функциональных тождества гамма-функции. Принципиально, что больше
ничего не требуется: все соотношения между используемыми
$\Ga$-произведениями выводятся из этих двух тождеств.

\begin{lemma}[формула отражения]\label{lem:reflection}
Для всех $z\in\CC\setminus\ZZ$
\[
  \Ga(z)\,\Ga(1-z)\;=\;\frac{\pi}{\sin \pi z}.
\]
\end{lemma}

\begin{proof}
Ключевое звено --- интеграл Эйлера
$\int_0^\infty\frac{u^{z-1}}{1+u}\,du=\frac{\pi}{\sin\pi z}$
при $0<\re z<1$, доказываемый вычетами по ключевому контуру:
функция $f(u)=\frac{(-u)^{z-1}}{1+u}$ голоморфна вне разреза
$[0,+\infty)$; единственный полюс $u=-1$ внутри контура имеет вычет
$1$; стороны разреза дают
$\bigl(e^{i\pi(1-z)}-e^{-i\pi(1-z)}\bigr)\int_0^\infty
\frac{x^{z-1}}{1+x}dx=2\pi i$. Замена $t=\frac{u}{1+u}$ переводит
интеграл в $\Bfun(z,1-z)=\Ga(z)\Ga(1-z)$; сопоставление даёт
формулу; аналитическое продолжение распространяет её на все
$z\notin\ZZ$.
\end{proof}

\begin{lemma}[формула умножения Гаусса]\label{lem:gauss}
Для целых $m\ge2$ и $mz\notin\ZZ_{\le0}$
\[
  \Ga(z)\,\Ga\!\Bigl(z+\frac1m\Bigr)\cdots
  \Ga\!\Bigl(z+\frac{m-1}{m}\Bigr)
  \;=\;(2\pi)^{\frac{m-1}{2}}\;m^{\frac12-mz}\;\Ga(mz).
\]
\end{lemma}

\begin{proof}
Отношение левой и правой частей $Q_m=G_m/H_m$ --- $1$-периодическая
голоморфная функция: рекуррентность $\Ga(w+1)=w\Ga(w)$ даёт
$\frac{G_m(z+1)}{G_m(z)}=\prod_{k=0}^{m-1}(z+\tfrac km)
=m^{-m}\frac{\Ga(mz+m)}{\Ga(mz)}$, и тот же множитель у $H_m$.
Формула Стирлинга даёт $Q_m(z)\to1$ при $\re z\to+\infty$ в
вертикальных полосах конечной ширины; периодичность переносит
предел на всю полосу, откуда $Q_m\equiv1$.
\end{proof}

\begin{lemma}[инвентарь соотношений]\label{lem:relations}
Все используемые соотношения между
$\Omega_{a,b}=\Bfun(a/N,b/N)$ --- следствия двух тождеств:
\begin{enumerate}[leftmargin=2em]
\item если $a+b=N$: $\Omega_{a,b}=\pi/\sin(\pi a/N)$ --- чистая
  дробная доля $\pi$;
\item если $a+b>N$:
  $\Omega_{a,b}=\frac{N}{a+b-N}
  \frac{\Ga(a/N)\Ga(b/N)}{\Ga((a+b-N)/N)}$ --- редукция;
\item если $a+b<N$: $\Omega_{a,b}$ канонично.
\end{enumerate}
\end{lemma}

\begin{proof}
Пункт 1 --- отражение при $b/N=1-a/N$; пункт 2 ---
рекуррентность $\Ga\bigl(\frac{a+b}N\bigr)=\frac{a+b-N}{N}
\Ga\bigl(\frac{a+b-N}N\bigr)$; пункт 3 --- по определению.
\end{proof}

\begin{databox}[числа]
Лестница отражения (V6): $\Ga(k/N)\Ga(1-k/N)=\pi/\sin(\pi k/N)$ для
всех $k=1,\dots,N-1$, $N=15,30$; максимум относительного отклонения
$7.0\cdot10^{-36}$ при $dps=35$.
\end{databox}

\begin{statusbox}[итог]
Оба тождества выведены из интегральных представлений и основной
теоремы о вычетах --- внешние результаты не привлекаются. Инвентарь
фиксирует: полная классификация $\QQ$-соотношений между
$\Ga$-произведениями (Рорлих) не используется --- достаточность
порождающего набора проверяется рангом (V7: ранг $=g$).
\end{statusbox}

\begin{verifybox}[как воспроизвести]
\texttt{python3 laboratory.py} $\to$ <<Стенд N=15/30 $\to$ V6>> ---
лестница отражения для всех $k$; диаграмма
\texttt{fig\_gamma.png} ($600$~dpi) --- численный портрет тождества.
\end{verifybox}

\vfill
{\color{linkblue}\hrule height 0.6pt}
\vspace{0.5em}
{\small\color{gray}
Лицензия: индивидуальная исключительная лицензия Исаева Исхака Хамзатовича.
Все права на вычисления, формулы и тексты принадлежат автору.
Верификация: \texttt{python3 laboratory.py} (пункт <<Стенды>>/<<Сертификаты>>),
Lean: \texttt{verification/lean/HodgeLaboratory.lean}.
}
\end{document}
